\input{preamble}

\title{Section 3.2 --- Root Properties --- Answer Key}

\begin{document}
\maketitle
\section*{Solving Equations with Root Properties}
Solve each equation using the appropriate root property. In this section, all roots will be integers. If there are multiple solutions, list them all. If there is no solution, write ``no solution.''
\begin{smenum}
\mitemxxxx
{$x=4$}
{$x=-5,5$}
{$x=-3,3$}
{$x=-2$}
\mitemxxxx
{$x=-9,9$}
{$x=-3$}
{No solution}
{$x=5$}
\mitemxxxx
{$x=-4,1$}
{$x=\frac{1}{3}$}
{$x=-\frac{3}{2},-1$}
{$x=1$}
\mitemxxxx
{$x=-8,-2$}
{$x=1,5$}
{$x=7$}
{$x=2$}
\end{smenum}
\section*{Simplifying Roots}
Simplify each root.
\begin{smenum}
\mitemxxxx
{$\sqrt{50}=5\sqrt{2}$}
{$\sqrt[3]{56}=2\sqrt[3]{7}$}
{$\sqrt{216}=6\sqrt{6}$}
{$\sqrt[3]{336}=2\sqrt[3]{42}$}
\end{smenum}
Simplify each fraction. If necessary, rationalize the denominator. Also simplify the radical in the numerator, if possible.
\begin{smenum}
\mitemxxxx
{$\sqrt{\dfrac{5}{16}}=\dfrac{\sqrt{5}}{4}$}
{$\sqrt[4]{\dfrac{11}{81}}=\dfrac{\sqrt[4]{11}}{4}$}
{$\sqrt[3]{\dfrac{25}{27}}=\dfrac{\sqrt[3]{25}}{3}$}
{$\sqrt{\dfrac{18}{49}}=\dfrac{3\sqrt{2}}{7}$}
\mitemxxxx
{$\sqrt[4]{\dfrac{7}{27}}=\dfrac{\sqrt[4]{21}}{3}$}
{$\sqrt[3]{\dfrac{49}{72}}=\dfrac{\sqrt[3]{147}}{6}$}
{$\sqrt{\dfrac{5}{200}}=\dfrac{\sqrt{10}}{20}$}
{$\sqrt[3]{\dfrac{125}{12}}=\dfrac{5\sqrt[3]{18}}{6}$}
\end{smenum}
\section*{Simplifying Solutions of Equations}
Solve each equation. Simplify your answer as much as possible. In this section, your answers may or may not be rational. If there are two solutions, list them both using $\pm$ notation and as two separate solutions.
\begin{smenum}
\mitemxxx
{$x=5\pm\sqrt{7}$\\$x=5+\sqrt{7},5-\sqrt{7}$}
{$x=3-\sqrt[3]{2}$}
{$x=4\pm\sqrt{2}$\\$x=4+\sqrt{2},4-\sqrt{2}$}
\mitemxxx
{$x=7+2\sqrt[3]{3}$}
{$x=\dfrac{15\pm 2\sqrt{15}}{3}$\\$x=\dfrac{15+2\sqrt{15}}{3},\dfrac{15-2\sqrt{15}}{3}$}
{$x=\dfrac{5+3\sqrt[3]{2}}{4}$}
\end{smenum}
\end{document}
